\section{Geometric Interpretation}

We interpret the transport construction geometrically, with particular
attention to the relationship between the cocycle and the sector
configuration on $\Gfifteen$.

\subsection{Two kinds of structure on the core}

The quotient graph $\Gfifteen \cong L(\mathrm{Petersen})$ carries two
distinct structures arising from transport:
\begin{itemize}
\item a $\mathbb{Z}_2$-valued cocycle $\sigfun$ on edges,
\item a family of $15$ transport sectors, one associated to each vertex.
\end{itemize}

The cocycle is sparse after switching reduction, while the sector system
is highly distributed.

\subsection{The cocycle support}

Let
\[
S = \{ e \in E(\Gfifteen) : \sigfun(e) = -1 \}.
\]

Before switching reduction, the negative-edge sector determined by the
transport signing has support on a connected $14$-edge subgraph. Its
vertex-incidence profile is nonuniform: five vertices have incidence
count $4$, eight vertices have incidence count $1$, and two vertices do
not occur at all.

This already shows that the cocycle support singles out a nontrivial
combinatorial region of the core.

\subsection{The sector geometry}

Independently, the transport sectors define a $15\times 30$ incidence
matrix $M$, whose centered row vectors produce a spherical
configuration. The centered inner products depend only on distance in
$\Gfifteen$ and take exactly three values:
\[
2.4666666667,\qquad -1.5333333333,\qquad -2.5333333333,
\]
corresponding respectively to distances
\[
1,\qquad 2,\qquad 3.
\]

After normalization, the corresponding cosines are
\[
0.3303571429,\qquad -0.2053571429,\qquad -0.3392857143.
\]

\subsection{Relation between the two}

The cocycle and the spherical sector geometry are related, but they do
not collapse to the same object.

The cocycle enters the construction through transport parity and the
signed lift, while the sector geometry is governed by the global
sector--edge overlap law
\[
MM^{\mathsf T} = A^3 + 2A^2 + 2I.
\]

In particular, the minimal cocycle support of size $6$ does not appear
to define a single eigenspace, orbit, or isolated geometric component of
the centered sector configuration. Rather, the cocycle contributes to
the transport structure from which the sectors arise, while the sector
geometry records a more distributed overlap pattern across the full
core.

\subsection{Interpretation}

The two structures should therefore be viewed as complementary:
\begin{itemize}
\item the cocycle captures a global $\mathbb{Z}_2$ obstruction in the
transport system,
\item the sector configuration captures the higher-order overlap geometry
induced by transport on the core.
\end{itemize}

Together they show that the transport system encodes more than the bare
graph structure of $\Gfifteen$.

\subsection{Relation to known graphs}

The graph $\Gsixty$ is a connected tetravalent graph on $60$ vertices
with diameter $6$, matching the standard identifier
\[
\mathrm{AT4val}[60,6].
\]

The significance of the present construction is not merely that it
recovers this graph, but that it equips it with an explicit signed-lift
structure and a rigid sector geometry descending to the Petersen core.
